% Chapter 5: Second to fourth derivatives (methods author). PLAN.md section 1.4; sidebars S8 and S9; figures F10
% and F11; examples E7 and E8 (E8's listing is quoted in Chapter 8 since round 1). No table (Chapter 9 has them).
% Revised at stage 6 per reviews/R1-decision.md (T9, T13, engineer-17, engineer-19, engineer-20, developer-16,
% library-12, evidence-21); the changes are listed in responses/R1-05-orders.md.
% Language pass for students (PLAN.md section 12.9): the wording only; content, evidence and order as at 3192e8a.
\chapter{Second to fourth derivatives}\label{ch:orders}

Many problems need more than the first derivative. The slope of a heat capacity, $\dd \cp / \dd T$, is the second
derivative of the enthalpy. The curvature that a solver needs to find a minimum is a second derivative too. A
stiffness or a Taylor coefficient may need the third or the fourth. The library computes derivatives of order 1 to 4
with every method, through one more template argument. This chapter shows the rules that it uses. It explains why
each order costs digits, and how many. Then it shows what the default step does for each order, what the float
warning means, and how Ridders' method extends to higher orders. The nitrogen example's second derivative,
$\enth''(T_0) = \dd \cp / \dd T$ at \SI{1000}{\kelvin}, runs through the chapter.

\section{The rules for N = 2 to 4}\label{sec:orders-rules}

The derivative order is the second template argument of \code{derivative}. It is either an integer or an enumerator
of \code{order}. So \code{derivative<Central<2>, 2>(fn, x)} and \code{derivative<Central<2>, order::second>(fn, x)}
are the same call. If you leave the order out, you get the first derivative: \code{derivative<M>} is
\code{derivative<M, 1>}. The rules behind the call come from the same moment conditions as in
\cref{sec:stencils-moments}, with $\mom{\dorder} = \dorder!$ in place of $\mom{1} = 1$. Take the second derivative
on the nodes $-1, 0, 1$. A symmetric rule has the weights $a, b, a$. The conditions are $\mom{0} = 2a + b = 0$ and
$\mom{2} = 2a = 2! = 2$. The second condition gives $a = 1$, and then the first gives $b = -2$. The odd moments
vanish by symmetry, since the nodes $-1$ and 1 carry the same weight; in particular $\mom{3} = 0$. So
\begin{equation}
  f''(x) = \frac{f(x - h) - 2 f(x) + f(x + h)}{h^2} - \frac{h^2}{12} f^{(4)}(\xi),
  \label{eq:orders-second}
\end{equation}
with the error from $\mom{4} = 2$ over $4!$, that is, $2/24 = 1/12$. This is the three-point rule. In the library it
is \code{Central<2>} for $\dorder = 2$: the same tag as for $f'$, with the same meaning. \code{Central<P>} is the
symmetric rule whose error is $\order(h^{\aorder})$, for every derivative order. For the $\dorder$-th derivative it
uses the $\aorder + \dorder - 1$ integers around 0. It includes node 0 for even $\dorder$ and leaves it out for odd
$\dorder$: by antisymmetry, an odd derivative gives the centre no weight (\cref{fig:stencil-nodes}). So
\code{Central<2>} evaluates $f$ three times for $f''$, four times for $f'''$ on the nodes $\pm 1, \pm 2$, and five
times for $f''''$:
\mnote{A rule that samples $x$ itself fails with \code{non_finite} where $f(x)$ is not finite. For example,
$\sin x / x$ at 0 works for $f'$ and $f'''$ with the central rules, but not for $f''$.}
\begin{align*}
  f'''(x)  &\approx \frac{f(x + 2h) - 2 f(x + h) + 2 f(x - h) - f(x - 2h)}{2 h^3}, \\
  f''''(x) &\approx \frac{f(x + 2h) - 4 f(x + h) + 6 f(x) - 4 f(x - h) + f(x - 2h)}{h^4},
\end{align*}
which are the centred rules of $\order(h^2)$ in Chapra's tables (Chapter 23).\cite{chapra2015} The $\order(h^4)$
rules of those tables are \code{Central<4>}. The one-sided families work for every order too. \code{Forward<P>} uses
the nodes $0, \ldots, \aorder + \dorder - 1$, and \code{Backward<P>} their mirror image. Their amplification grows
quickly with $\dorder$ (\cref{tab:weights-forward-fourth}), but exactly at an edge they are what works. One detail:
the library evaluates its three-point rule for $f''$ in a slightly different form from \cref{eq:orders-second}. It
uses the offsets of the points as they were actually rounded (\cref{sb:actual-offsets}). The two forms agree in exact
arithmetic whenever the points are exact. The complete tables for $\dorder = 2$ to 4 are in \cref{ch:tables}.

\section{Why accuracy falls with the order}\label{sec:orders-accuracy}

\begin{figure}
\sidecaption{The V-shaped error curve of \code{Central<2>} for $\dorder = 1$ to 4, for $\mathrm{e}^x$ at 1. The
truncation branches nearly coincide. The rounding branch steepens as $1/h^{\dorder}$, and the minimum rises. This is
why the fourth derivative gives about \num{\valHeadlineFourthDigits} digits where the first gives
\num{\valHeadlineFirstDigits}.\label{fig:v-curve-orders}}
\centering
% F10 fig:v-curve-orders (Chapter 5, text width). The V of Central<2> for N = 1 to 4, exp(x) at x = 1 (DS5,
% data/step-error-orders.csv). The truncation branches coincide (every derivative of exp is exp, and the constants
% 1/6, 1/12, 1/4, 1/6 differ little); the rounding branch steepens as 1/h^N and the minimum rises with N: the fourth
% derivative (accent) gives about four digits where the first gives eleven. Errors above 1 are cut off.
\begin{tikzpicture}
\begin{axis}[eqtufte, width=95mm, height=45mm,
  xmode=log, ymode=log, log basis x=10, log basis y=10,
  xmin=1e-13, xmax=1, ymin=1e-13, ymax=1,
  xtick={1e-12,1e-10,1e-8,1e-6,1e-4,1e-2,1}, ytick={1e-12,1e-9,1e-6,1e-3,1},
  xlabel={step $h$}, ylabel={relative error}, restrict y to domain=-14:-0.05, unbounded coords=jump]
  \addplot[black!80] table[col sep=comma, comment chars={\#}, x=h, y=n1] {data/step-error-orders.csv};
  \addplot[black!65] table[col sep=comma, comment chars={\#}, x=h, y=n2] {data/step-error-orders.csv};
  \addplot[black!50] table[col sep=comma, comment chars={\#}, x=h, y=n3] {data/step-error-orders.csv};
  \addplot[accent] table[col sep=comma, comment chars={\#}, x=h, y=n4] {data/step-error-orders.csv};
  % labels below and to the right of the rounding branches, where nothing else is drawn
  \node[font=\scriptsize, anchor=north west, inner sep=1pt] at (axis cs:3e-10, 1e-9) {$\dorder = 1$};
  \node[font=\scriptsize, anchor=north west, inner sep=1pt] at (axis cs:3e-7, 1e-5) {$\dorder = 2$};
  \node[font=\scriptsize, anchor=east, inner sep=1pt] at (axis cs:1.2e-5, 3e-2) {$\dorder = 3$};
  \node[accent, font=\scriptsize, anchor=south west, inner sep=1pt] at (axis cs:3e-4, 8e-2) {$\dorder = 4$};
\end{axis}
\end{tikzpicture}

\end{figure}

Each derivative order costs digits, and this section shows why. \Cref{fig:v-curve-orders} is the V of
\cref{ch:rounding}, drawn for each derivative order. Again the function is $\mathrm{e}^x$ at $x = 1$, so that the
orders compare. On the right, the truncation branches nearly coincide. There are two reasons: every derivative of
$\mathrm{e}^x$ is $\mathrm{e}^x$, and the error constants of the four rules, $\tfrac16$, $\tfrac1{12}$, $\tfrac14$
and $\tfrac16$, differ little. On the left, the rounding branches do not coincide. A rounding error of
$\epsmach \abs{f}$ in each sample is divided by $h^{\dorder}$. So on logarithmic axes the branch has slope
$-\dorder$. Halving the step doubles the rounding error of $f'$, and multiplies that of $f''''$ by sixteen, which is
$2^4$. The error model of \cref{sec:stencils-price} becomes
\begin{equation}
  E(h) = \errc\, h^{\aorder} \abs{f^{(\aorder + \dorder)}} + \ampl\, \epsmach \frac{\abs{f}}{h^{\dorder}},
  \label{eq:orders-model}
\end{equation}
and its minimum, worked out in \cref{sec:derivations-model}, is of order $\epsmach^{\aorder/(\aorder + \dorder)}$.
For \code{Central<2>} the accuracy order is $\aorder = 2$, so the exponent is $2/(2 + \dorder)$. In double that gives
$\epsmach^{2/3}$, $\epsmach^{1/2}$, $\epsmach^{2/5}$ and $\epsmach^{1/3}$ for $\dorder = 1$ to 4. These are
\num{\valModelCentralTwoFirstDouble}, \num{\valModelCentralTwoSecondDouble}, \num{\valModelCentralTwoThirdDouble}
and \num{\valModelCentralTwoFourthDouble}, or about \num{\valModelDigitsCentralTwoFirstDouble},
\num{\valModelDigitsCentralTwoSecondDouble}, \num{\valModelDigitsCentralTwoThirdDouble} and
\num{\valModelDigitsCentralTwoFourthDouble} digits. The minima of \cref{fig:v-curve-orders} rise accordingly. The
measured medians at the default step, in \cref{fig:order-accuracy}, follow the model within its usual factor. They
are \num{\valDoubleFourthCentralRelError} for the fourth derivative, against \num{\valExpCentralTwoRelError} for the
first. These are the \num{\valHeadlineFourthDigits} and \num{\valHeadlineFirstDigits} digits that \cref{ch:why}
promised. \Cref{sb:fourth-sensitivity} quantifies what is lost to each order, and why float cannot afford a
derivative above the first.

One caution about the word \emph{relative}. The model's rounding term, $\ampl \epsmach \abs{f} / h^{\dorder}$, is an
absolute error. To turn it into a relative error, we divide it by $\abs{f^{(\dorder)}}$. At the default step this
gives the relative error $\ampl \epsmach^{\aorder/(\aorder + \dorder)}$ times the factor
$\abs{f} / (\scale^{\dorder} \abs{f^{(\dorder)}})$. For $\mathrm{e}^x$ at 1 that factor is 1. The figures of
\cref{fig:order-accuracy} presuppose that $f$ varies on the scale $\scale$, so that $\abs{f^{(\dorder)}}$ is about
$\abs{f} / \scale^{\dorder}$ and the factor is about 1. For the nitrogen example's second derivative at $T_0$, the
factor $\abs{\enth} / (\scale^2 \abs{\enth''})$ is \num{\valCpSecondRelativeFactor}. So the model's relative error is
$\ampl = 4$ (the sum of the absolute weights $1, 2, 1$) times that factor times $\epsmach^{1/2}$, which is
\num{\valCpSecondModelRelError}. The measured error of \num{\valCpSecondRelError} in \cref{sec:orders-thread} lies
below it. Near a point where $f^{(\dorder)}$ vanishes, the relative error grows without bound, while the absolute
error stays what the model says. There, judge a result by its absolute error (\cref{sec:using-check}).

\begin{sidebar}[label=sb:fourth-sensitivity]{Why the fourth derivative is so sensitive}
By \cref{eq:orders-model}, the rounding error of a rule for $f^{(\dorder)}$ is
$\ampl \epsmach \abs{f} / h^{\dorder}$. So each halving of the step multiplies it by $2^{\dorder}$. As a result, the
balance with truncation is reached at a larger step, and at a larger error. At the best step the error is of order
$\epsmach^{\aorder/(\aorder + \dorder)}$. With \code{Central<2>}, the exponent falls from $\tfrac23$ for $f'$ to
$\tfrac13$ for $f''''$. In double, that is the difference between about \num{\valModelDigitsCentralTwoFirstDouble}
digits and about \num{\valModelDigitsCentralTwoFourthDouble}. \code{Central<4>} does better, from $\epsmach^{4/5}$ to
$\epsmach^{1/2}$, at the cost of two more evaluations (\cref{tab:default-steps}).

In float, $\epsmach$ is \num{\valEpsFloat}, and the same exponents leave much less. Here $\epsmach^{2/3}$ is
\num{\valModelCentralTwoFirstFloat} for $f'$, and $\epsmach^{1/3}$ is \num{\valModelCentralTwoFourthFloat} for
$f''''$: two digits, by the model. Measured on $\mathrm{e}^x$, the median error of \code{Central<2>} for the fourth
derivative in float is \num{\valFloatFourthCentralRelError}. Ridders reaches \num{\valFloatFourthRiddersRelError},
with a relative estimate of \num{\valFloatFourthRiddersRelEstimate}, so the estimate reports the noise honestly. The
second and third derivatives keep more: \num{\valFloatSecondCentralRelError} and \num{\valFloatThirdCentralRelError}
with \code{Central<2>}. But that is far less than the \num{\valFloatFirstCentralRelError} of $f'$. This is why every
derivative above the first carries the warning of \cref{sec:orders-float}.

The sensitivity to the step grows the same way. The rounding term is proportional to $h^{-\dorder}$, and the
truncation term to $h^{\aorder}$. So a step a factor 2 below the right one multiplies the rounding term by
$2^{\dorder}$. A step a factor 2 above it multiplies the truncation term by $2^{\aorder}$. The default step assumes
that $f$ varies on the scale $\scale$. Suppose that we start at the model's balance and misjudge that scale by a
factor 2. Then the total error grows by at most \num{\valMisjudgeCostCentralTwoFirst} for $f'$ and
\num{\valMisjudgeCostCentralTwoFourth} for $f''''$ with \code{Central<2>}, and by at most
\num{\valMisjudgeCostCentralFourFourth} for \code{Central<4>} and $f''''$. The library's documentation quotes a
factor 16 for $f''''$. That is the growth of the rounding term alone, $2^{\dorder}$. A step on $f$'s own scale
(\cref{sec:rounding-scale}) removes the sensitivity. Ridders removes the need to know the scale at all.
\end{sidebar}

\begin{figure}
\sidecaption{The relative error at the default step against $\dorder$, for \code{Central<2>}, \code{Central<4>} and
Ridders, in double and in float, on $\mathrm{e}^x$ (medians over eleven points). The dotted line for each type is
the model $\epsmach^{\aorder/(\aorder + \dorder)}$ of \code{Central<2>}, as a guide. The float points for $f''$ to
$f''''$ are boxed: these are the cases that warn.\label{fig:order-accuracy}}
\centering
% F11 fig:order-accuracy (Chapter 5, text width). Relative error at the default step against the derivative order N
% for Central<2>, Central<4> and Ridders, in double (solid) and float (dashed), exp(x), medians over eleven points
% in [0.5, 1] (DS6, data/order-accuracy.csv). The dotted guides are the model eps^(2/(2+N)) of Central<2> for each
% type (2^-52 and 2^-23: exact constants). The float fourth derivatives are boxed in accent: the warning case.
\begin{tikzpicture}
\begin{axis}[eqtufte, width=76mm, height=48mm,
  ymode=log, log basis y=10, xmin=0.7, xmax=4.3, ymin=1e-15, ymax=1,
  xtick={1,2,3,4}, ytick={1e-15,1e-12,1e-9,1e-6,1e-3,1},
  xlabel={derivative order $\dorder$}, ylabel={relative error at the default step}]
  % the model of Central<2>: eps^(P/(P+N)) with P = 2, for double and float (dotted; named in the caption)
  \addplot[muted, dotted, domain=1:4, samples=25] {(2^(-52))^(2/(2+x))};
  \addplot[muted, dotted, domain=1:4, samples=25] {(2^(-23))^(2/(2+x))};
  % double, solid
  \addplot[black!80, mark=*, mark size=1.2pt] table[col sep=comma, comment chars={\#}, x=N, y=c2_double] {data/order-accuracy.csv};
  \addplot[black!80, mark=square*, mark size=1.2pt] table[col sep=comma, comment chars={\#}, x=N, y=c4_double] {data/order-accuracy.csv};
  \addplot[black!80, mark=triangle*, mark size=1.5pt] table[col sep=comma, comment chars={\#}, x=N, y=rid_double] {data/order-accuracy.csv};
  % float, dashed, the same marks
  \addplot[black!60, dashed, mark=*, mark size=1.2pt, mark options={solid}] table[col sep=comma, comment chars={\#}, x=N, y=c2_float] {data/order-accuracy.csv};
  \addplot[black!60, dashed, mark=square*, mark size=1.2pt, mark options={solid}] table[col sep=comma, comment chars={\#}, x=N, y=c4_float] {data/order-accuracy.csv};
  \addplot[black!60, dashed, mark=triangle*, mark size=1.5pt, mark options={solid}] table[col sep=comma, comment chars={\#}, x=N, y=rid_float] {data/order-accuracy.csv};
  % labels at the right end of the double curves; the float curves carry the same marks
  \node[font=\scriptsize, anchor=west, inner sep=1pt] at (axis cs:4.07, 1.2e-5) {\code{Central<2>}, double};
  \node[font=\scriptsize, anchor=west, inner sep=1pt] at (axis cs:4.07, 4.9e-8) {\code{Central<4>}, double};
  \node[font=\scriptsize, anchor=west, inner sep=1pt] at (axis cs:4.07, 6.8e-10) {\code{Ridders}, double};
  \node[font=\scriptsize, anchor=south, inner sep=1pt] at (axis cs:2.5, 4e-3) {float, dashed (same marks)};
  % the warning case: every fourth derivative in float
  \draw[accent, thin] (axis cs:3.85, 3e-4) rectangle (axis cs:4.15, 4e-2);
  \node[accent, font=\scriptsize, anchor=south] at (axis cs:4, 4.5e-2) {float $f''''$: warning};
\end{axis}
\end{tikzpicture}

\end{figure}

\section{The default step per order and its sensitivity}\label{sec:orders-step}

How does the default step change with the order? It balances the two terms of \cref{eq:orders-model}, under the
assumption that $f$ varies on the scale $\scale = \max(\abs{x}, 1)$. Under that assumption,
$\abs{f^{(\aorder + \dorder)}}$ is about $\abs{f} / \scale^{\aorder + \dorder}$. Setting the two terms equal, with the
constant dropped as in \cref{sb:step-constant}, gives $h^{\aorder + \dorder} = \epsmach \scale^{\aorder + \dorder}$,
that is,
\begin{equation}
  h = \epsmach^{1/(\aorder + \dorder)}\, \scale .
  \label{eq:orders-step}
\end{equation}
The exponent is $1/(\aorder + \dorder)$, so the step grows with the order: the exponent shrinks, and since
$\epsmach$ is far below 1, a smaller exponent gives a larger number. For \code{Central<2>} in double the step is
$\epsmach^{1/3}$, $\epsmach^{1/4}$, $\epsmach^{1/5}$ and $\epsmach^{1/6}$ times $\scale$. At $\scale = 1$ that is
\num{\valRootEpsThreeDouble}, \num{\valRootEpsFourDouble}, \num{\valRootEpsFiveDouble} and
\num{\valRootEpsSixDouble}. At $T_0 = \SI{1000}{\kelvin}$, the step for $\enth''$ is \SI{\valCpSecondStep}{\kelvin},
against \SI{\valCpDefaultStep}{\kelvin} for $\enth'$. \Cref{tab:default-steps} lists every rule's exponent, step and
model error.

The constant that the model would add is $(\dorder \ampl / (\aorder \errc))^{1/(\aorder + \dorder)}$. With
\code{Central<2>} it is \num{\valStepConstantCentralTwoSecond} for $f''$ and \num{\valStepConstantCentralTwoFourth}
for $f''''$. By the model, dropping it costs more as $\dorder$ grows. Yet on test functions that vary faster than
$\scale$, half the default step was better for every order (\cref{sb:step-constant}). So the library keeps one
convention, and documents the sensitivity instead (\cref{sec:derivations-model}).

That sensitivity is the practical point. Take a step that is off by a factor 2 from the right one. By the model, the
total error grows by up to \num{\valMisjudgeCostCentralTwoFirst} for $f'$, and by up to
\num{\valMisjudgeCostCentralTwoFourth} for the fourth derivative, with \code{Central<2>}
(\cref{sb:fourth-sensitivity}). The rounding term alone grows by $2^{\dorder}$. When $f$ varies on a scale $w$ much
smaller than $\scale$, pass $h = \epsmach^{1/(\aorder + \dorder)} w$, as \cref{sec:rounding-scale} does for $f'$. And
the default step depends on both the method and the order. So generic code that varies $\dorder$ should call the
default-step overload, rather than pass one $h$ for every order.

\section{The float warning}\label{sec:orders-float}

In float, every derivative above the first keeps few correct digits. Here are the median errors, measured on
$\mathrm{e}^x$. \code{Central<2>} is \num{\valFloatFirstCentralRelError} off for $f'$, but
\num{\valFloatSecondCentralRelError}, \num{\valFloatThirdCentralRelError} and \num{\valFloatFourthCentralRelError}
for $f''$ to $f''''$. In double it is \num{\valDoubleFourthCentralRelError} for $f''''$ (\cref{fig:order-accuracy},
\cref{sb:fourth-sensitivity}). Ridders does better: from \num{\valFloatSecondRiddersRelError} for $f''$ to
\num{\valFloatFourthRiddersRelError} for $f''''$. Its one-sided form of \cref{sec:ridders-onesided} does worse: from
\num{\valFloatSecondForwardRiddersRelError} for $f''$ to \num{\valFloatFourthForwardRiddersRelError} for $f''''$,
which is more than the value itself.

The library tells you so at compile time. Take a call for a derivative above the first, by any method, in a type
whose epsilon is above $10^{-7}$. That is float, decimal32 or bfloat16, but not double. Such a call selects a
deprecated overload of \code{derivative},\mnote{A \emph{deprecated} function still works, but the compiler warns at
every call of it.} and the compiler reports at the caller's line: ``derivatives above the first in a type with 7 or
fewer significant digits keep few correct digits: use double, or silence this warning at the call or with
NUMERIXX\_NO\_PRECISION\_WARNING''. The cure is to use double. If you want a float result all the same, you can
silence the warning for one call or for the whole build. The result is the same with or without the warning
(\cref{sec:api-warning}, which quotes the call and its numbers). Ridders' estimate then says how little the result
is worth: \num{\valFloatSecondRiddersRelEstimate} relative for $f''$ and \num{\valFloatFourthRiddersRelEstimate} for
$f''''$.

\section{Ridders for higher orders}\label{sec:orders-ridders}

Ridders' method extends to every order with the same tableau.\mnote{Ridders' method, the subject of the next
chapter, computes a base rule at a sequence of halved steps, the \emph{levels}, and extrapolates them in a triangular
table, the \emph{tableau}.} The base of level $i$ is the \code{Central<2>} rule for $\dorder$, at the step
$h_0 / 2^i$. Its error expansion has even powers of $h$ only, so the extrapolation formula of \cref{ch:ridders}
applies unchanged. Three things differ. First, the default initial step is $\scale / 4$ for $\dorder \ge 2$, rather
than $\scale / 32$. The base's rounding error grows as $1/h^{\dorder}$, so a larger start pays. Second, a level costs
2 evaluations for $f'$ and $f''$, and 4 for $f'''$ and $f''''$. The even orders also need $f(x)$, once per call, and
it is kept across restarts. Third, the gain over a stencil grows with the order. For the fourth derivative in double,
Ridders reaches a median of \num{\valDoubleFourthRiddersRelError}, where \code{Central<2>} reaches
\num{\valDoubleFourthCentralRelError}. That is four more digits, for a median of \num{\valRiddersEvalsMedianFourth}
evaluations over the random study of \cref{ch:ridders}, against five. \Cref{fig:order-accuracy} shows the three
methods side by side.

\section{The nitrogen example: $\dd \cp / \dd T$}\label{sec:orders-thread}

Now back to nitrogen. By the closed form, the slope of the heat capacity at $T_0$ is
$\enth''(T_0) = \SI{\valCpSecondExact}{\joule\per\mole\per\kelvin\squared}$. Two methods compute it, in one call
each: the three-point rule with its default step of \SI{\valCpSecondStep}{\kelvin}, and Ridders from
$h_0 = T_0 / 4$:
\snippet{05-orders-second}{second-derivative}
The stencil's relative error is \num{\valCpSecondRelError}, below the model's \num{\valCpSecondModelRelError} of
\cref{sec:orders-accuracy}. Ridders returns a value \num{\valCpRiddersSecondRelError} off. Its estimate is
\SI{\valCpRiddersSecondEstimate}{\joule\per\mole\per\kelvin\squared}, against an actual error of
\SI{\valCpRiddersSecondActual}{\joule\per\mole\per\kelvin\squared}, after \num{\valCpRiddersSecondEvals}
evaluations. Finally, the integer and the enumerator spell the same order. The two calls below return equal results,
value and estimate alike:
\snippet{05-orders-second}{order-enum}

\begin{sidebar}[label=sb:actual-offsets]{The three-point $f''$ from the actual offsets}
The two-node rules of \cref{sb:actual-distance} divide by the distance between the points as evaluated. This
removes the rounding of the points. The library's three-point rule for $f''$ does the same with its two offsets. Let
$x_+ = \fl(x + h)$ and $x_- = \fl(x - h)$ be the two outer points as rounded. The rule forms the offsets
$h_+ = x_+ - x$ and $h_- = x - x_-$, and evaluates
\[
  f''(x) \approx \frac{\dfrac{f(x_+) - f(x)}{h_+} - \dfrac{f(x) - f(x_-)}{h_-}}{(h_+ + h_-) / 2},
\]
that is, the second divided difference of the three points where they actually are. When $h_+ = h_- = h$, it is the
three-point rule again.

The offsets $h_+$ and $h_-$ are exact for every $0 < h \le \abs{x}$, in any radix, by Sterbenz's lemma
(Theorem 4.3.1).\cite{sterbenz1974} The lemma says that the difference of two floating-point numbers of the same sign,
within a factor 2 of each other, is computed exactly. There are two cases. For $0 < h \le \abs{x} / 2$, $x_+$ and
$x$ lie within a factor 2 of each other and have the same sign, and so do $x$ and $x_-$. So the subtractions round
nothing. For $\abs{x} / 2 < h \le \abs{x}$, the subtraction $x - h$ is itself exact by the same lemma. So
$x_- = x - h$, and the offset $h_-$ is $h$, a representable number, with $x_- = 0$ at $h = \abs{x}$. Meanwhile $x_+$
stays within a factor 2 of $x$, as before.

With exact offsets, the formula is exact for every quadratic, whatever the rounding of the points. The point-rounding
term that a nominal-step formula carries, about $\epsmach \abs{x f'} / h^2$, drops out. This mattered most for
Ridders in float, whose second-derivative levels use the same form. In the prototype of the design record
(\code{docs/deriv/PHASE4_API.md}), the nominal step gave a worst case of 8.7 times the estimate, and the actual
offsets 1.4. The form also never computes $h^2$, which underflows in double for $h$ below about $10^{-154}$. For the
same reason, every rule with more nodes divides by $\den h$ and then by $h$ again, $\dorder - 1$ times. It does so in
the order that avoids a spurious underflow or overflow: by $h$ first when $\den h > 1 > h$. Only \code{Central<2>} for
$f''$ uses the offsets. The rules with more nodes divide by the nominal step, as \cref{sec:stencils-price} explains.
Ridders' levels for $f'''$ and $f''''$ instead account for the point rounding in their floor
(\cref{sb:floor-coverage}).
\end{sidebar}
